\documentclass[11pt, a4paper]{article}

\usepackage[utf8]{inputenc}
\usepackage[T1]{fontenc}
\usepackage[a4paper, margin=2.5cm]{geometry}
\usepackage{enumitem}
\setlist[itemize]{label=-}
\usepackage{hyperref}
\usepackage{graphicx}
\usepackage{array}
\usepackage{booktabs}

\title{\textbf{Plataformas de Monitorização de QoS e QoE na Internet: Análise da Infraestrutura CAIDA Ark}}
\author{
  \textbf{GRUPO 6} \\[0.4cm]
  Ivo Costa Sousa -- \texttt{pg63976} \\[0.15cm]
  João Afonso Almeida Sousa -- \texttt{pg62405} \\[0.15cm]
  João Carlos Teixeira Neiva -- \texttt{pg63978} \\[0.15cm]
}
\date{Outubro de 2026}

\begin{document}

\maketitle

\begin{abstract}
\textit{Este relatório, desenvolvido para a Unidade Curricular de Qualidade de Serviço na Internet, foca-se na análise da infraestrutura CAIDA Ark. O estudo explora o seu impacto prático no QoS e QoE, identificando limitações técnicas das sondagens ativas e propondo melhorias realistas para a monitorização de redes.}
\end{abstract}

\section{Parte 1: Pesquisa e Resumo do CAIDA Ark}

\subsection{Propósito}
O \textbf{CAIDA Ark} realiza medições ativas à escala global para mapear a topologia da Internet e detetar congestionamentos interdomínio. Ao contrário de ferramentas comerciais para o consumidor final, o sistema destina-se a investigadores e operadores de rede que necessitam de analisar o comportamento global e a estabilidade estrutural da Internet.

\subsection{Métricas Chave Medidas}
A plataforma assenta na execução coordenada de sondagens ativas, recorrendo a ferramentas especializadas (como o \textit{scamper}), recolhendo os seguintes indicadores principais:
\begin{itemize}
    \item \textbf{Caminhos de Roteamento (IP Paths / Topology):} A sequência exata de saltos (roteadores e Sistemas Autónomos) que um pacote de dados atravessa desde a origem até ao destino.
    \item \textbf{Round-Trip Time (RTT):} O tempo de latência de ida e volta, medido em milissegundos, para cada salto individual ao longo da rota examinada.
    \item \textbf{Perda de Pacotes (Packet Loss):} A percentagem de pacotes de teste que são descartados ou não recebem resposta ao longo do trajeto.
\end{itemize}

\subsection{Relação com QoS e QoE}
A recolha de dados do CAIDA Ark estabelece uma ponte direta com os parâmetros de avaliação de rede:
\begin{itemize}
    \item \textbf{Mapeamento para QoS:} Os dados traduzem-se diretamente em medidas básicas de Qualidade de Serviço. O RTT estabelece a latência de ponta a ponta e o \textit{jitter} (através da variação do RTT ao longo do tempo). A topologia e a perda de pacotes são cruciais para localizar gargalos físicos ou saturação em links BGP nas bordas das redes.
    \item \textbf{Mapeamento para QoE:} Embora o CAIDA Ark não realize medições ao nível da aplicação de forma direta, as suas métricas de QoS servem de base para a perceção do utilizador final. Aumentos acentuados no RTT ou estrangulamentos identificados em ligações interdomínio causam problemas severos de \textit{lag} e dessincronização de \textit{hit registation} em jogos competitivos. De igual modo, conduzem a tempos de carregamento de páginas web mais longos ou a \textit{buffering} indesejado em serviços de \textit{streaming} de vídeo.
\end{itemize}

\section{Parte 2: Análise e Discussão}

\subsection{Comparação de Abordagens de Medição}
A medição de desempenho de redes de computadores pode ser efetuada através de diferentes paradigmas, cada um com as suas vantagens e utilidades distintas (ver Tabela \ref{tab:comparacao}).

As abordagens ativas, como os testes efetuados com \textit{Speedtest} ou as sondagens contínuas do \textit{CAIDA Ark}, baseiam-se na injeção de tráfego artificial na rede para aferir a capacidade máxima, a latência e mapear rotas. Estas métricas são muito uteis para identificar estrangulamentos de QoS ao nível da infraestrutura pura. Em contrapartida, as abordagens passivas, tipificadas pela plataforma \textit{SamKnows}, observam o tráfego real gerado pelo utilizador sem introduzir sobrecarga artificial na rede. Este método proporciona uma visão mais fiel da QoE, dado que reflete as condições reais que os utilizadores enfrentam no seu consumo diário.

No cenário específico do \textit{CAIDA Ark} e do \textit{MANIC}, o modelo de medição é estritamente ativo, sustentado no envio de pacotes coordenados de sondagem. Embora tenha uma excelente precisão na deteção de congestionamentos interdomínio (QoS), esta abordagem carece de triangulação com outras fontes de dados para traduzir o impacto real na experiência do utilizador (QoE), visto que a simples topologia e o estado do protocolo BGP não espelham o desempenho na camada aplicacional.

\begin{table}[hbt!]
\centering
\caption{Comparação entre Abordagens de Medição (Ativa vs. Passiva)}
\label{tab:comparacao}
\begin{tabular}{@{} >{\raggedright\arraybackslash}p{2.8cm} >{\raggedright\arraybackslash}p{4.5cm} >{\raggedright\arraybackslash}p{4.5cm} @{}}
\toprule
\textbf{Parâmetro} & \textbf{Ativa (Ex: CAIDA Ark, Speedtest)} & \textbf{Passiva (Ex: SamKnows)} \\
\midrule
\textbf{Mecanismo} & Injeção de tráfego artificial/teste. & Observação e escuta de tráfego real. \\
\addlinespace
\textbf{Foco Principal} & QoS: limites da infraestrutura, rotas, latência máxima. & QoE: experiência efetiva da aplicação e do utilizador final. \\
\addlinespace
\textbf{Impacto na Rede} & Pode gerar sobrecarga (\textit{overhead}) temporária nos \textit{links}. & Impacto nulo ou residual na capacidade de transmissão. \\
\addlinespace
\textbf{Ponto Forte} & Isolamento preciso de problemas estruturais e estrangulamentos. & Fidelidade absoluta às condições de utilização diária e contínua. \\
\bottomrule
\end{tabular}
\end{table}

\subsection{Contribuição para a Melhoria da Internet}
Plataformas abertas como o CAIDA Ark, o M-Lab e o RIPE Atlas desempenham um papel importante na transparência da Internet. A \textit{open source} ajuda ativamente na gestão, regulação e otimização global da rede. A partilha destes \textit{datasets} possibilita que a comunidade científica e os órgãos reguladores verifiquem a neutralidade da rede e arbitrem de forma isenta disputas de \textit{peering} entre diferentes ISPs.

Adicionalmente, a identificação sistemática e documentada de congestionamentos recorrentes em ligações de trânsito força os operadores de telecomunicações a justificarem as suas políticas e a negociarem reforços de capacidade, tais como a adição de novas portas de 100G/400G em pontos de troca de tráfego, com base em evidências inegáveis. No domínio puramente académico, os dados longitudinais do CAIDA têm sido um pilar na modelação do crescimento do ecossistema IPv6 e no estudo da resiliência global a falhas ou ataques BGP \textit{hijacking}.

\subsection{Limitações e Viés}
Qualquer sistema de monitorização à escala global está sujeito a limitações. A vulnerabilidade mais notória do CAIDA Ark reside no seu viés geográfico. Estando dependente do alojamento voluntário e institucional de monitores, existe uma concentração esmagadora de sondas na América do Norte e na Europa, deixando muitas regiões de África e da Ásia com representação estatística marginal. Esta assimetria conduz a uma visão fraturada da topologia global e impossibilita uma avaliação equitativa da QoS em economias emergentes.

Ao nível técnico, as medições ativas deparam-se frequentemente com a filtragem intencional e a despriorização de pacotes ICMP efetuada pelos \textit{routers} na rede de \textit{core}. Como demonstrado na análise do nosso grupo ao fenómeno dos \textit{Orphan Flows}, esta interceção gera falsos positivos nas estatísticas de perda de pacotes e distorce o cálculo do \textit{jitter}. Para mitigar estes erros, torna-se essencial a aplicação de modelos estatísticos de regressão/compensação, paralelamente a um esforço diplomático para estabelecer parcerias com ISPs locais nas geografias sub-representadas para a instalação de novos pontos de presença.

\subsection{Melhorias no Monitoramento}
Para contornar as falhas de cobertura e as atuais limitações das medições ativas, sugerimos a integração do \textit{CAIDA Ark} com redes de monitorização mais abrangentes, como o RIPE Atlas. Ao utilizar os milhares de pontos de recolha de dados deste projeto, conseguiríamos melhorar significativamente o mapeamento da topologia global da Internet. Desta forma, o módulo MANIC passaria a dispor de informação mais realista e exata para conseguir detetar os congestionamentos entre os diferentes domínios de rede.

Uma alternativa viável para melhorar o sistema seria a adoção do \textit{eBPF} (Extended Berkeley Packet Filter). Como opera diretamente no kernel do Linux, esta tecnologia permite recolher dados da rede com uma precisão muito elevada e com baixo consumo de processamento. Consideramos que esta abordagem tornaria as medições de QoS muito mais fiáveis, uma vez que reduz a dependência de pacotes de teste,que são frequentemente bloqueados, e ajuda a detetar variações mínimas na latência.

\subsection{Tendências Futuras}
O futuro da monitorização de redes exigirá sistemas descentralizados baseados em \textit{Edge Computing} e modelos de \textit{Machine Learning}. Para dar resposta a infraestruturas cada vez mais exigentes, acreditamos que plataformas como o CAIDA Ark terão de transitar da sua atual recolha periódica para uma abordagem preditiva. Dessa forma, ao processar localmente as variações de RTT, os monitores conseguirão antecipar congestionamentos em rotas BGP antes que estes afetem a experiência do utilizador, tornando a gestão da rede totalmente proativa.

\end{document}